\documentclass[referee]{aa}

\usepackage{graphicx}
\usepackage{txfonts}
\usepackage{natbib}
\usepackage{hyperref}
\usepackage{caption}
\usepackage{lipsum}
\usepackage{subcaption}         
\usepackage{lscape}             
\usepackage{placeins}           
\usepackage{booktabs}
\usepackage{amsmath}

\begin{document}

    \title{Determining the long-term dynamical and temperature stability of potentially habitable planets in trinary star
    systems GJ 667 and LTT 1445}


   \titlerunning{Stability of exoplanets in hierarchical triple systems}
    \authorrunning{Donev}
    
    \author{Ema Donev\inst{1}}

    \institute{
    University of Groningen, PO Box 72, 9700 AB Groningen, The Netherlands\\
    \email{emadonev@icloud.com}
    }
    
    \date{Received XXX; accepted XXX}
    
    \abstract
      {}
      {Investigate long-term dynamical and temperature stability on exoplanets of systems Gliese 667, LTT 1445, HD 132563, Kepler-444, and 94 Ceti.}
      {Used a Wisdom-Holman integrator adapted for triple hierarchical star systems and temperature tracking. For systems without orbital parameters for stellar orbits, explored a range of eccentricities and inclinations.}
      {2 out of 2 habitable exoplanets in Gliese 667 remain stable long term, 94 Ceti and Kepler 444 support the existence of synthetic exoplanets, while LTT 1445 is highly dynamically unstable and HD 132563 has very high temperature changes.}
      {With this work long-term stability can be determined, but not necessarily habitability by only using the habitable zone and temperature estimates.}
    
    \keywords{ exoplanets -- N-body simulation -- habitability
    }
    
    \maketitle

    \section{Introduction} \label{sec:intro}
    The search for potentially habitable planets has always been a core part of astrophysics. 
    From the first exoplanet discovered in 1992 around pulsar PSR~B1257+12 \citep{WolszczanFrail1992} to today, 
    more than 6000 exoplanets have been listed in the NASA Exoplanet Archive \citep{NASAExoplanetArchiveCounts2026}.
    A non-negligible fraction of known exoplanets orbit stars that are members of binary or higher-order multiple systems,
    although the inferred fraction depends on the sample and observational methods \citep[e.g.][]{Roell2012, 
    MichelMugrauer2021,Schlagenhauf2024, Raghavad2006}. Due to the complex orbital architecture of multiple star systems
    and their exoplanets, analyzing the dynamical stability of such systems is crucial for long-term habitability. The presence of 
    other stellar companions can be the driver of the Kozai-Lidov mechanism \citep[][]{Kozai1962, Lidov1962}, influencing the temperature and climate
    of a potentially habitable exoplanet. Stable zones in binary systems have been widely studied, with boundaries of stable zones
    empirically defined for all types of planetary orbits \citep{HolmanWiegert1999}. Related work studied habitable zones of 
    binary systems \citep[e.g.][]{HaghighipourKaltenegger2013, KalteneggerHaghighipour2013, EgglGeorgakarakos2020}. For triple hierarchical systems most
    body of work focuses on dynamical stability \citep[][]{VerrierEvans2007,Busetti2018}, with more recent work exploring habitable zones in such systems
    \citep[e.g.][]{Muller2014}. 
    Even though there are only 73 triple star systems listed in the NASA Exoplanet archive \citep{NASAExoplanetArchiveCounts2026}
    , they offer a wide range of planetary configurations and potential habitable exoplanets. 2 very interesting
    and close cases, within 24 light years, are systems Gliese 667 and LTT 1445. Gliese 667 C hosts 2 Super-Earths, with GJ 667 Cc being in the habitable zone and an exciting
    candidate for habitability
    \citep[]{AngladaEscudéGJ667, DelfosseHARPS}, although there has been significant debate over the total number of planets. 
    \citep{Anglada2013multipleplanets} claimed there are 6, or even 7 planets in the system with 3 of them being in the habitable zone, while later
    research \citep[][]{FerozHobsonBayesian2014, Robertson2014} argues that only GJ 667 Cb and GJ 667 Cc are confirmed, with the other planets being artifacts
    on the RV periodograms. While GJ 667 Cc appears to be a prime target for habitability, its closeness to its host M dwarf causes multiple issues affecting habitability, such as
    rising temperatures due to spin-down time \citep{MarakovGJ667C} and flares of UV and X-rays which can harm any potential life if the planet has no magnetic field \citep{SloaneGJ667cXray}.
    LTT 1445 A hosts 2 planets as well, LTT 1445 Ab and LTT 1445 Ac with a possible third candidate, LTT 1445 Ad. \citep[][]{Winters_2019, Winters_2022}
    Planets LTT 1445 Ac and Ad have potential to be habitable, as LTT 1445 Ac has a similar radius and density as the Earth \citep{Pass_2023}, making it an exciting candidate for further studies,
    as well as LTT 1445 Ad which is in the habitable zone. However, LTT 1445 Ab is also an interesting candidate for further investigation of its atmosphere, as \citep{Bennett_2025} found traces of HCN and a $~2\sigma$ certainty
    of an atmosphere. The rocky planets in this system also face similar issues as in GJ 667, with X rays from the host star LTT 1445 A having the potential to harm
    water retention, although planet LTT 1445 Ad could still keep its water supply if it contains one. %\citep{keylist}.
    In this work the goal is to investigate the long-term dynamical stability of planets GJ 667 Cb and Cc, and LTT 1445 Ab, Ac and Ad, determine the influence of the trinary system on the dynamics of the planets
    and the effective temperature, as well as track how the temperature changes as the planets progress in their dynamical evolution using N-body integrators. The data used is presented in
    \hyperref[sec:data]{Sect. 2}, the methodology of an adapted N-body integrator in \hyperref[sec:method]{Sect. 3}, the results and discussion in \hyperref[sec:result]{Sect. 4} and \hyperref[sec:discuss]{Sect. 5},
    and finally the conclusions of this work in \hyperref[sec:conclusion]{Sect. 6}.

    \section{Data}\label{sec:data}

    To investigate the long-term dynamical and temperature stability of potentially habitable exoplanets, orbital and photometric data from all stellar and planetary components
    is required. Tables \ref{tab:stellar}, \ref{tab:orbits} and \ref{tab:planets} contain
    all the data used in this work. The luminosities of the stellar components will be used to calculate the effective surface temperature
    on the exoplanets, further explained in section \ref{sec:temp}. To ensure all the orbital data is in the same reference plane and the planets start at the correct phase, 
    a conversion is introduced from various $T_0$ to the true anomaly or $\theta$. In the binary system B-C in LTT 1445, $T_0$ in Besselian years is converted to
    the true anomaly by calculating the mean anomaly and using a universal Kepler solver to calculate the value of $\theta$. A similar procedure is done for the outer orbit in LTT 1445,
    however the mean anomaly is $M = \lambda_{ref} - \omega$. The planets in LTT 1445 have a time of conjunction, so a conversion to the time of periastron is needed first, and then
    the procedure follows the previous examples. The following expressions were used for this conversion:
    \begin{equation}
        \theta_{conj} = \frac{\pi}{2} - \omega \\
        E_{conj} = 2\arctan[\sqrt{\frac{1-e}{1+e}} \cdot \tan(\theta_{conj}/2)] \\
        M_{conj} = E_{conj} - e\sin E_{conj} \\
        T_0 = T - \frac{P}{2\pi} \cdot M_{conj}
    \end{equation}
    Finally, for the LTT 1445 planets, $\Omega$ will be set to 0, while the inclunation for planet d (as it is non-transiting) will be set to the same inclination
    as the other two planets under the assumption of coplanarity at the angle of $89^\circ$. The N-body code will be fed data from the sky-view, as precise orbits of all the planets and stellar
    components are available and measured as sky inclinations.
    Since the stellar components of GJ 667 are poorly constrained, the binary pair AB will be merged into a single object and the orbit of C - (A-B) will be investigated through
    a study of the parameter space. Noted with $(S)$ in Table \ref{tab:orbits}, the range of values
    for eccentricity is \{$0.0, 0.2, 0.4, 0.6, 0.8$\} and for mutual inclination \{$0^\circ, 30^\circ, 60^\circ, 90^\circ$\} (with this selection, two values are below the critical
    Kozai-Lidov angle of $39.2^\circ$, and two above it). The values of the remaining orbital elements are $\Omega = 0^\circ$, $\omega = 0^\circ$ and $\theta_0 = 90^\circ$.
    As the confirmed GJ 667 planets are non-transiting, coplanarity is assumed with a mutual inclination of 0 and $\Omega = 0$. Mutual inclination is used in this system as there is no
    measured sky inclination quantity so the coordinate system is built on mutual inclinations of orbits AB around C and the planets' orbits.

\begin{table*}[t]
        \centering
        \small
        \caption{Stellar properties of the trinary sytems GJ 667 and LTT 1445.}
        \label{tab:stellar}
        \begin{tabular}{lccccccccc}
                \toprule
                \textbf{Star} &
                \multicolumn{4}{c}{\textbf{LTT 1445}} &
                \multicolumn{4}{c}{\textbf{GJ 667}} \\
                \cmidrule(lr){2-5}
                \cmidrule(lr){6-9} 
                & $M (M_{\odot})$ & $R (R_{\odot})$ & $T (K)$ & $L_{bol}(L_{\odot})$ &
                $M (M_{\odot})$ & $R (R_{\odot})$ & $T (K)$ & $ L_{bol} (L_{\odot})$ \\
                \midrule
                A & $0.257 \pm 0.014$ & $0.268 \pm 0.027$ & $3337 \pm 150$ & $0.00794^{(D)}$ & - & $0.76$ & $4270$ & $0.173^{(D)}$ \\
                B  & $0.215 \pm 0.014$ & $0.236 \pm 0.027$ & - & $0.00596^{(D)}$ & - & $0.70$ &  $3998$ & $0.113^{(D)}$ \\
                C  & $0.161 \pm 0.014$ & $0.197 \pm 0.027$ & - & $0.00368^{(D)}$ & $0.327 \pm 0.008$ & $0.337 \pm 0.014$  &  ${3443^{+75}_{-71}}$ & $0.01439 \pm 0.00034$\\
                \bottomrule
        \end{tabular}
        \tablefoot{$(D)$ denotes a value calculated or unit-converted. The luminosities of LTT 1445 A,B and C from \citep{Rukdee2024} are converted in terms of the Sun's luminosity according to \citep{Prsa2016}. 
        The masses, radii and temperatures of LTT 1445 stellar components are adopted from \citet{Winters_2019, Rukdee2024}. The masses of components A and B in the GJ 667 system
        currently have no explicit values, however according to \citet{Desgrange2023} the AB pair has a combined mass of $1.27 M_{\odot}$ and can be treated
        as a single object for N-body integration. Luminosities of stars GJ 667 A and B are adopted Stefan-Boltzmann values, $L/L_{\odot} = (R/R_{\odot})^2 (T_{eff}/5772 K)^4$, calculated from older component-resolved radii and temperatures 
        \citet{JohnsonWright1983}. The properties of the host, GJ 667 C are adopted from \citet{Pineda2021}.}

    \end{table*}


    \begin{table*}[t]
        \centering
        \small
        \caption{Orbital properties of the trinary sytems GJ 667 and LTT 1445.}
        \label{tab:orbits}
        \begin{tabular}{lccccc}
                \toprule
                \textbf{} &
                \multicolumn{2}{c}{\textbf{LTT 1445}} &
                \multicolumn{1}{c}{\textbf{GJ 667}} \\
                \cmidrule(lr){2-3}
                \cmidrule(lr){4-4} 
                & $B - C$ & $A - (B - C)$ &
                $C - (A - B)$ & \\
                \midrule
                $P \; $ (yr)& $36.2 \pm 5.3$ & ${549^{+243}_{136}}$  & -\\
                $\alpha \; (\arcsec)$  & $1.159 \pm 0.076$ & - & $32.4$ \\
                $a \; $ (au) & $8.1 \pm 0.5$ & ${58^{+16}_{-10}}$ & $~235$ \\
                $e $  & $0.50 \pm 0.11$ & ${0.375^{+0.084}_{-0.037}}$ &  $[0.0, 0.2, 0.4, 0.6, 0.8]^{(S)}$\\
                $i_{sky} \; (\deg)$  & $89.64 \pm 0.13$ & ${88.5^{+1.3}_{-1.4}}$ & - \\
                $i_{mut} \; (\deg)$  & - & $2.88 \pm 0.63$ & $[0.0, 30.0, 60.0, 90.0]^{(S)}$ \\
                $\Omega \; (\deg) $  & $137.63 \pm 0.19$ & $135.15 \pm 0.28$ & - \\
                $\omega \; (\deg) $  & $209 \pm 13$ & ${256^{+28}_{-21}}$ & - \\
                $T_0 \; $ (Besselian yr)  & $2019.2 \pm 1.7$ & - & - \\
                $\lambda_{ref} (\deg)$  & - & $210.3 \pm 2.1$ & - \\
                \bottomrule
        \end{tabular}
        \tablefoot{Orbital data for the binary pair A-B in LTT 1445 was adopted from \citet{Winters_2019}, while the data for the outer pair
        is adopted from \citet{DYnamicalZhangLTT1445}. The orbital parameters listed for the outer orbit of LTT 1445 in the table are the posterior median values,
        but in the evolutions of planetary systems the best-fit values will be used. Since there is no precise estimate of the separate masses of components A and B in GJ 667,
        their orbit is irrelevant as they will be approximated as a single body. The orbit of component C around the merged binary AB doesn't have a fully defined orbit \citep{Desgrange2023},
        so a range of values for eccentricity and inclination will be included as testing the parameter space for possible stable orbits of this system, denoted by $(S)$.}

    \end{table*}

    \begin{table*}[t]
        \centering
        \small
        \caption{Properties of exoplanets in trinary sytems GJ 667 and LTT 1445.}
        \label{tab:planets}
        \begin{tabular}{lcccccc}
                \toprule
                \textbf{} &
                \multicolumn{3}{c}{\textbf{LTT 1445}} &
                \multicolumn{2}{c}{\textbf{GJ 667}} \\
                \cmidrule(lr){2-4}
                \cmidrule(lr){5-6} 
                & b & c & d &
                b & c & \\
                \midrule
                $M (M_\oplus)$ & $2.32 \pm 0.25$ & $1.00 \pm 0.19$ & - & - & - &\\
                $M \sin i (M_\oplus)$ & - & - & $2.72 \pm 0.75$ & $5.6 \pm 0.3$ & $4.1 \pm 0.6$ &\\
                $R (R_\oplus)$  & $1.43 \pm 0.09$ & ${1.07^{+0.10}_{-0.07}}$ & - & - & - &\\
                $a \; (au)$  & $0.03807$ & $0.02656$ & $0.09$ & $0.050431 \pm 0.000004$ & $0.12501 \pm 0.00009$ &\\
                $e $  & ${0.18^{+0.07}_{-0.06}}$ & ${0.05^{+0.05}_{-0.03}}$ &  ${0.26^{+0.14}_{-0.13}}$& $0.15 \pm 0.05$ & $0.27 \pm 0.1$ &\\
                $i_{sky} \; (\deg)$  & $88.95^{+0.64}_{-0.57}$ & ${86.5^{+0.39}_{-0.47}}$ & - & - & - &\\
                $i_{mut} \; (\deg)$  & - & - & - & - & - &\\
                $\Omega \; (\deg) $  & - & - & - & - & - & \\
                $\omega \; (\deg) $  & ${243.5^{+14.9}_{-17.8}}^{(D)}$ & ${162.7^{+10.3}_{-108.3}}^{(D)}$ & ${50.4^{+49.8}_{-33.8}}^{(D)}$ & $12 \pm 20$ & $140 \pm 20$ &\\
                $T_0 \;$(BJD)  & ${2458423.426378^{+0.000299}_{-0.000304}}$ & ${2457005.358765 \pm 0.000003}$ & ${2457000.04618^{+0.00085}_{-0.00086}}$ & $2454446.1 \pm 0.3$ & $2454462 \pm 2$ & \\
                \bottomrule
        \end{tabular}
        \tablefoot{Orbital information and the parameters of planets b, c and d are adapted from \citet{Lavie2023}, with the radius
        of planet c from \citet{Pass_2023} due to higher precision. The values of orbital parameters are taken as eccentric posteriors from \citet{Lavie2023}, as 
        eccentricity is a key component in habitability. For GJ 667 the values are adopted from \citet{Robertson2014}. }

    \end{table*}

    \section{Methodology}\label{sec:method}

    - exploring dynamical and temperature stability

    - guideline for LTT 1445 is Verier et all 2007 and the main architecture

    - for GJ 667 due to the outer component being a single object, use binary code system and main source

    - link code and github

    \subsection{Temperature calculation and climate models}\label{sec:temp}
    - reference and explanation for temperature calculation

    - explain methodology of code and the testing for every system
    
    \section{Results}\label{sec:result}
    
    \section{Discussion}\label{sec:discuss}
    
    \section{Conclusions}\label{sec:conclusion}
    
    \begin{acknowledgements}
    Write funding, telescope/archive acknowledgements, software acknowledgements, and thanks here.
    This research has made use of the NASA Exoplanet Archive, which is operated
    by the California Institute of Technology, under contract with the National
    Aeronautics and Space Administration under the Exoplanet Exploration Program.
    \end{acknowledgements}
    
    \bibliographystyle{aa}
    \bibliography{reference}
    
    \begin{appendix}
    \section{Additional material}
    Appendix text here.
    \end{appendix}
    
\end{document}